\section*{Chapitre \ref{ch:rc:market-risk-measures} --- Mesures du risque de marché}

\begin{solution}{exo:rc:market-risk-measures:1}
$k = \lceil0.01\times500\rceil = 5$~: la cinquième plus forte perte est la VaR. Pour l'ES à 97.5~\%,
$\lceil0.025\times500\rceil = 13$~: la moyenne des treize plus fortes pertes.
\end{solution}

\begin{solution}{exo:rc:market-risk-measures:2}
$1.58\text{ million}\times\sqrt{10} = \text{USD}~5.01$ million. La règle échoue lorsque les variations sont
autocorrélées (tendances, retour à la moyenne), lorsque la volatilité évolue sur les dix jours, lorsque les
positions ne sont pas conservées (couverture dynamique) et pour les portefeuilles non linéaires, dont le P\&L sur dix jours
n'est pas égal à dix P\&L à un jour.
\end{solution}

\begin{solution}{exo:rc:market-risk-measures:3}
Zone jaune, avec une hausse de 0.65~: un facteur de multiplication de 3.65.
\end{solution}

\begin{solution}{exo:rc:market-risk-measures:4}
Pris isolément~: $\P(L>0) = 0.9~\% < 1~\%$, donc la VaR vaut 0~; le pire 1~\% contient le défaut de 0.9~\% (100) et 0.1~\%
de zéros, ES $= 0.9\times100/1 = 90$. Ensemble~: $\P(L>0) = 1-0.991^2 = 1.79~\% > 1~\%$ tandis que $\P(L>100) =
0.0081~\% < 1~\%$, donc la VaR vaut 100~; le pire 1~\% contient 0.0081~\% de 200 et 0.9919~\% de 100, ES $= 100.8$.
\end{solution}

\begin{solution}{exo:rc:market-risk-measures:5}
Un straddle vendeur a un gamma négatif~: il perd $\tfrac12\Gamma(\Delta S)^2$ sur des mouvements
dans un sens comme dans l'autre. Une réévaluation en delta seul omet cette perte, la plus forte
précisément dans les scénarios de queue qui fixent la VaR~; la VaR historique vaut USD~1.91 million avec le seul delta et 2.10 avec le
terme en gamma.
\end{solution}

\begin{solution}{exo:rc:market-risk-measures:6}
La volatilité de suivi fixe bien l'échelle en moyenne, mais la loi normale place trop
peu de probabilité dans les queues~: au quantile de 99~\%, les mouvements à queues épaisses dépassent $2.33\hat\sigma$ plus
souvent que dans 1~\% des cas. Les dépassements viennent de la forme, non du niveau.
\end{solution}

\begin{solution}{exo:rc:market-risk-measures:7}
\texttt{kupiec(1000, 28, 0.01)} donne un rapport de vraisemblance de 22.0 et une p-valeur de
$3\times10^{-6}$~: le modèle est rejeté.
\end{solution}

\begin{solution}{exo:rc:market-risk-measures:8}
Aucun dépassement en 250 jours survient avec une probabilité de 8.1~\% pour un modèle correct~: un faible indice de
conservatisme. Sur quatre ans, le modèle historique a connu exactement 1~\% de dépassements, tandis que le modèle EWMA,
bien que son taux soit plus bas, a échoué au \omterm{def:rc:market-risk-measures:tests}{test de Kupiec} avec un taux de 2.8~\%. Plus bas ne signifie pas meilleur.
\end{solution}

\begin{solution}{pb:rc:market-risk-measures:1}
\textbf{1.} EWMA~: VaR USD~1.58 million, ES 1.59 million~; historique~: 2.10 et 2.34 millions.
\textbf{2.} L'exposition au rendement à dix ans~: 114~\% de la VaR EWMA, partiellement compensée par la position courte à deux ans.
\textbf{3.} USD~5.01 million.
\textbf{4.} Sous normalité, l'ES à 97.5~\% vaut 2.338 écarts-types et la VaR à 99~\% 2.326.
\textbf{5.} Les chiffres historiques, ou les deux avec une explication~: le conseil a besoin du
chiffre à queues épaisses et de son incertitude, non du plus bas.
\textbf{6.} 8 dépassements, zone jaune.
\textbf{7.} $3+0.75 = 3.75$.
\textbf{8.} USD~15.02 million à 3, 18.78 million à 3.75, en prenant la VaR à dix jours d'aujourd'hui pour la moyenne sur soixante jours.
\textbf{9.} Aucun.
\textbf{10.} Oui~: avec 3 dépassements, il serait en zone verte sans hausse, bien que sa VaR soit plus élevée
(USD~2.09 million)~; la comparaison des fonds propres dépend du niveau face au multiplicateur.
\textbf{11.} Sa forme est inadaptée~: la queue de la loi normale est trop fine pour ces facteurs.
\textbf{12.} L'\omterm{def:rc:market-risk-measures:es}{expected shortfall} à 97.5~\% sur un horizon de base de dix jours mis à l'échelle par les horizons de liquidité,
avec un backtest au niveau du desk et le test d'attribution du P\&L (chapitre~23).
\textbf{13.} 8 dépassements à 99~\% restent dans les 12 autorisés~; le décompte à 97.5~\% devrait aussi rester
sous 30.
\textbf{14.} Savoir si les dépassements se regroupent~: un modèle peut avoir le bon nombre mais échouer par rafales,
ce qui est pire pour l'adéquation des fonds propres.
\textbf{15.} \omterm{def:rc:market-risk-measures:fhs}{Simulation historique filtrée}~: queues épaisses et volatilité rapide, en zone verte sur quatre
ans~; sa VaR est plus élevée, mais sans hausse du multiplicateur ses fonds propres restent comparables.
\textbf{16.} Le rôle d'un modèle est de bien décrire la queue~; un chiffre bas, dépassé trop souvent,
sous-estime les fonds propres et trompe sur les limites de risque.
\textbf{17.} Oui~: en concentrant le risque dans des positions dont les pertes se situent au-delà du quantile (options
vendues, crédit), en négociant sur des fenêtres d'historique périmées, ou en exploitant des facteurs que le modèle omet.
\textbf{18.} L'ES, les résultats de stress, les plus fortes contributions, la concentration et la liquidité, la
consommation des limites, et les résultats de backtest.
\textbf{19.} Résultat nommé~: \emph{le rapport de quatre heures et quart}, une VaR à 99~\% de USD~1.58 million et une
ES à 97.5~\% de 1.59 million, 8 dépassements en 250 jours, et un multiplicateur de 3.75.
\textbf{20.} Une perte qui n'est dépassée qu'avec une faible probabilité, dans les conditions récentes~; et non
l'ampleur des pertes au-delà, ni ce qui se passe lorsque les conditions changent.
\end{solution}

\begin{solution}{iq:rc:market-risk-measures:1}
VaR~: le quantile de perte à un niveau de confiance donné sur un horizon. ES~: la perte moyenne au-delà de ce
quantile. L'ES voit la taille de la queue et est cohérente (sous-additive)~: on ne peut la contourner en
dissimulant du risque au-delà du quantile, et elle ne pénalise pas la diversification.
\iqlookfor{les définitions et les arguments sur la queue et la cohérence.}
\end{solution}

\begin{solution}{iq:rc:market-risk-measures:2}
Historique~: mouvements conjoints réels et queues épaisses, réévaluation complète, mais réaction lente et limitée à
la fenêtre historique. Paramétrique~: rapide et analytique, linéaire et normale. Monte Carlo~: toute distribution et
réévaluation complète, mais dépendante du modèle et coûteuse.
\iqlookfor{compromis sur les hypothèses, la rapidité et la non-linéarité.}
\end{solution}

\begin{solution}{iq:rc:market-risk-measures:3}
Deux obligations indépendantes, chacune en défaut avec une probabilité de 0.9~\%~: chacune a une VaR à 99~\% nulle, le couple une VaR de 100.
\iqlookfor{un contre-exemple concret.}
\end{solution}

\begin{solution}{iq:rc:market-risk-measures:4}
Compter les jours dont la perte dépasse la VaR de la veille~; le \omterm{def:rc:market-risk-measures:tests}{test de Kupiec} vérifie si la fréquence correspond
au niveau choisi~; le \omterm{def:rc:market-risk-measures:tests}{test de Christoffersen} vérifie si les dépassements sont indépendants dans le temps~; le \omterm{def:rc:market-risk-measures:traffic}{test
des feux tricolores} associe les décomptes à des pénalités de fonds propres.
\iqlookfor{fréquence et indépendance.}
\end{solution}

\begin{solution}{iq:rc:market-risk-measures:5}
Réévaluer par type de transaction avec une valorisation vectorisée ou sur grille, des sensibilités pour les transactions linéaires, mettre en cache
les termes indépendants des scénarios, paralléliser par transaction et par scénario, exécuter de façon incrémentale
pour les transactions modifiées, et conserver la réévaluation complète pour la queue non linéaire.
\iqlookfor{grilles, sensibilités et parallélisme.}
\end{solution}

\begin{solution}{iq:rc:market-risk-measures:6}
L'ES n'est pas élicitable seule, mais les tests conjoints VaR--ES fonctionnent~: comparer les pertes de queue réalisées
à l'ES prédite les jours de dépassement, ou effectuer un backtest de la VaR à plusieurs niveaux
(le FRTB utilise 97.5~\% et 99~\%).
\iqlookfor{comparaisons des pertes de queue et tests de VaR à plusieurs niveaux.}
\end{solution}
